% !TEX program = pdflatex
% Hugging Face Deep RL Course - Bonus Unit 2 (full Russian translation).
% Compile with latexmk -pdf from notes/ru/ using ../preamble.tex.
\PassOptionsToPackage{T2A}{fontenc}
\PassOptionsToPackage{russian}{babel}
\documentclass[11pt,a4paper]{article}

% Shared typography/layout for Hugging Face Deep RL Course notes.
% Compile with pdfLaTeX / latexmk -pdf with notes/ as the working directory.
% e.g. cd notes && latexmk -pdf unit1.tex

\usepackage[T1]{fontenc}
\usepackage[utf8]{inputenc}
\usepackage[english]{babel}
\usepackage{lmodern}
\usepackage{textcomp}
\usepackage{geometry}
\geometry{a4paper,margin=22mm,headheight=14pt,headsep=8mm,footskip=12mm}

\usepackage{amsmath,amssymb,mathtools}
\usepackage{graphicx,caption,needspace,placeins}
\usepackage{array,booktabs,tabularx}
\usepackage{xcolor}
\usepackage{microtype}
\usepackage{enumitem}
\usepackage{listings}
\usepackage[most]{tcolorbox}
\usepackage{titlesec}
\usepackage{fancyhdr}
\usepackage{hyperref}
\usepackage{bookmark}

\definecolor{ink}{HTML}{243041}
\definecolor{accent}{HTML}{38598C}
\definecolor{secondary}{HTML}{60768A}
\definecolor{codebg}{HTML}{F4F6F8}
\definecolor{codeborder}{HTML}{E3E8EE}
\definecolor{notebg}{HTML}{F3F7FC}

\hypersetup{
  colorlinks=true,
  linkcolor=accent,
  urlcolor=accent,
  citecolor=accent,
  pdfborder={0 0 0}
}
\urlstyle{same}

\setlength{\parindent}{0pt}
\setlength{\parskip}{0.38em plus 0.12em minus 0.05em}

\setlist[itemize]{
  topsep=0.35em,
  itemsep=0.13em,
  leftmargin=1.7em
}
\setlist[enumerate]{
  topsep=0.35em,
  itemsep=0.13em,
  leftmargin=1.8em
}

\setcounter{tocdepth}{2}
\setcounter{secnumdepth}{2}

\titleformat{\section}
  {\Large\bfseries\color{ink}}
  {\thesection}{0.7em}{}

\titleformat{\subsection}
  {\large\bfseries\color{accent}}
  {\thesubsection}{0.7em}{}

\titleformat{\subsubsection}
  {\normalsize\bfseries\color{ink}}
  {\thesubsubsection}{0.6em}{}

\titlespacing*{\section}
  {0pt}{2.2ex plus 0.7ex}{0.9ex}
\titlespacing*{\subsection}
  {0pt}{1.7ex plus 0.4ex}{0.6ex}
\titlespacing*{\subsubsection}
  {0pt}{1.3ex plus 0.3ex}{0.45ex}

\pagestyle{fancy}
\fancyhf{}

\fancyhead[L]{
  \footnotesize\sffamily\color{secondary}
  HF Deep RL Course \enspace/\enspace Unit 1
}
\fancyhead[R]{
  \footnotesize\sffamily\color{secondary}
  \nouppercase{\rightmark}
}
\fancyfoot[C]{
  \footnotesize\color{secondary}\thepage
}
\renewcommand{\headrulewidth}{0.25pt}

\captionsetup{
  font=small,
  labelfont={bf,color=secondary},
  textfont={color=ink},
  justification=centering,
  skip=4pt,
  hypcap=false
}

\lstdefinestyle{hfcode}{
  basicstyle=\ttfamily\footnotesize,
  backgroundcolor=\color{codebg},
  frame=single,
  rulecolor=\color{codeborder},
  numbers=none,
  showstringspaces=false,
  breaklines=true,
  breakatwhitespace=false,
  columns=fullflexible,
  keepspaces=true,
  xleftmargin=0pt,
  xrightmargin=0pt,
  aboveskip=0.6em,
  belowskip=0.6em,
  framexleftmargin=7pt,
  framexrightmargin=7pt,
  framextopmargin=5pt,
  framexbottommargin=5pt,
  upquote=true,
  tabsize=4,
}
\lstset{style=hfcode}

\newcommand{\CourseFigure}[2]{%
  \par\addvspace{0.55em}\noindent%
  \begin{minipage}{\linewidth}%
    \centering
    \includegraphics[
      width=0.86\linewidth,
      height=0.295\textheight,
      keepaspectratio
    ]{#1}\par
    \captionof{figure}{#2}%
  \end{minipage}%
  \par\addvspace{0.55em}%
}

\newtcolorbox{coursenote}{
  enhanced,
  breakable,
  colback=notebg,
  colframe=accent!30,
  boxrule=0.35pt,
  arc=1.2mm,
  left=9pt,
  right=9pt,
  top=6pt,
  bottom=6pt
}

\renewenvironment{quote}
  {\begin{coursenote}}
  {\end{coursenote}}

\emergencystretch=2.5em

\providecommand{\tightlist}{
  \setlength{\itemsep}{0pt}
  \setlength{\parskip}{0pt}
}

\providecommand{\passthrough}[1]{#1}

\Urlmuskip=0mu plus 6mu\relax

\renewcommand{\rmdefault}{cmr}
\renewcommand{\sfdefault}{cmss}
\renewcommand{\ttdefault}{cmtt}
\graphicspath{{../assets/unitbonus2/}}
\fancyhead[L]{\footnotesize\sffamily\color{secondary}Курс Hugging Face Deep RL / Бонусный модуль 2}
\fancyhead[R]{\footnotesize\sffamily\color{secondary}Подбор гиперпараметров}
\title{\textbf{Глубокое обучение с подкреплением}\\[0.35em]\large Бонусный модуль 2: автоматический подбор гиперпараметров с Optuna}
\author{Курс Hugging Face Deep RL\\\small Полная учебная версия на основе MDX}
\date{}

\begin{document}
\selectlanguage{russian}
\fontencoding{T2A}\selectfont
\maketitle
\vspace{-1em}
\begin{center}\small\color{secondary}Источник: курс Hugging Face Deep RL (оригинальный текст на английском языке).\end{center}
\vspace{-0.45em}
\begingroup\small\setlength{\parskip}{0pt}\tableofcontents\endgroup
\vspace{0.3em}

\section{Введение}\label{introduction}

Одна из важнейших задач в глубоком обучении с подкреплением --- \textbf{найти удачный набор гиперпараметров обучения}.

\CourseFigure{optuna-logo.png}{Логотип библиотеки Optuna. \protect\href{https://raw.githubusercontent.com/optuna/optuna/master/docs/image/optuna-logo.png}{[оригинал изображения]}}

\href{https://optuna.org/}{Optuna} --- библиотека, которая помогает автоматизировать этот поиск. В этом модуле мы немного изучим \textbf{теорию автоматического подбора гиперпараметров}. Сначала попробуем вручную оптимизировать параметры алгоритма DQN, который рассматривали в предыдущем модуле. Затем \textbf{научимся автоматизировать поиск с помощью Optuna}.

\section{Руководство по Optuna}\label{optuna}

Материал ниже основан на \href{https://araffin.github.io/tools-for-robotic-rl-icra2022/}{презентациях Антонена Раффена (Antonin Raffin) для конференции ICRA 2022}. Он является одним из основателей проектов Stable-Baselines и RL-Baselines3-Zoo.

\subsection{Теория подбора гиперпараметров}

\begin{coursenote}
  \textbf{Видеолекция: теория подбора гиперпараметров}\par
  \href{https://www.youtube.com/watch?v=AidFTOdGNFQ}{Смотреть на YouTube (AidFTOdGNFQ)}.
\end{coursenote}

\subsection{Практическое руководство по Optuna}

\begin{coursenote}
  \textbf{Видеоурок: Optuna}\par
  \href{https://www.youtube.com/watch?v=ihP7E76KGOI}{Смотреть на YouTube (ihP7E76KGOI)}.
\end{coursenote}

Ноутбук $\rightarrow$ \href{https://colab.research.google.com/github/araffin/tools-for-robotic-rl-icra2022/blob/main/notebooks/optuna_lab.ipynb}{по этой ссылке}.

\section{Практические задания}\label{hands-on}

Теперь, когда ты научился пользоваться Optuna, вот несколько идей, как применить полученные знания:

\begin{enumerate}[label=\textbf{\arabic*.},leftmargin=1.8em,itemsep=0.6em]
  \item \textbf{Улучши результаты своего агента для LunarLander-v2}, используя Optuna для поиска более удачного набора гиперпараметров. Можешь также попробовать другую среду, например MountainCar-v0 или CartPole-v1.
  \item \textbf{Улучши результаты своего агента для SpaceInvaders}.
\end{enumerate}

Выполнив эти задания, ты увидишь, насколько полезной и мощной может быть Optuna при обучении более эффективных агентов.

Удачи и получай удовольствие от экспериментов!

И наконец, нам было бы очень интересно \textbf{узнать, что ты думаешь о курсе и как мы могли бы его улучшить}. Если у тебя есть отзывы или предложения, пожалуйста, $\rightarrow$ \href{https://forms.gle/BzKXWzLAGZESGNaE9}{заполни эту форму}.

\subsection*{Продолжай учиться и оставайся молодцом!}

\end{document}
